\section{Green Agent 设计清单：Task、Environment、Metrics、Test Cases}

slides 82--85 给出从目标到测试的四步。顺序很重要：先选要评价的 task，再设计 participant 运行环境，然后定义 metrics，最后用 test cases 覆盖成功、失败与边界。若先写 Agent 再找指标，通常会得到只验证 happy path 的项目。

\subsection{Step 1：选择明确且有价值的任务}

任务描述应包含用户目标、Agent 可用信息、允许动作、终止条件和不可违反约束。例如 ticket-booking agent 不只是“订票”，还包括日期/预算/偏好、库存变化、支付权限和取消规则。一个 benchmark 可以含多个 task family，但每类都要有清楚 outcome。

\lecturefigure{slides-images/slide-082.png}{官方 slide 82：选择要评价的 task}

\begin{warningbox}{避免把产品名称当任务定义}
“测试客服 Agent”过于宽泛；“在给定订单状态和政策下完成退款，保持审计字段一致且不越权”才可验证。任务越模糊，rubric 越依赖 judge 主观解释，模型也越容易用语言表面掩盖失败。
\end{warningbox}

\subsection{Step 2：环境必须可复现、隔离且表达真实约束}

任务目标确定后，环境决定 Agent 实际能看到和改变什么。本小节从 observation、action、permission 与 reset 四个接口检查复现性：环境既要保留真实约束，又必须隔离外部副作用，并能在每次 run 前回到可验证的初始状态。
slide 83 要求设计 Agent 需要运行的 environment：可交互 UI、数据库、API、工具和隐藏状态。环境初始化要确定，执行后可重置，外部依赖应 mock 或固定版本。若使用真实网络服务，要记录响应或接受不可复现性，并为故障与限流设计测试。

\lecturefigure{slides-images/slide-083.png}{官方 slide 83：环境、网页/API、动作与状态反馈}

\begin{knowledgebox}{Environment contract}
定义 observation schema、action schema、permission、timeout、resource budget、state transition 和 terminal condition。Green Agent 只通过 contract 控制环境，White Agent 只获得允许接口。每次 run 保存 initial/final state digest 和 tool trace，便于复现。
\end{knowledgebox}

\subsection{Step 3：Metrics 要覆盖成功、质量、成本与风险}

slide 84 以订票为例：不仅看是否订到，还看价格、行程是否满足用户要求。通用 metrics 可分 outcome success、constraint satisfaction、trajectory efficiency、robustness 和 safety。多指标可以形成 Pareto frontier；若必须汇总，权重和硬门槛必须公开。

\lecturefigure{slides-images/slide-084.png}{官方 slide 84：Green Agent 的 success rate 与用户约束指标}

\begin{importantbox}{硬门槛与软评分分离}
越权、数据泄露、环境破坏等应作为 hard fail，不应被高质量文字补偿；价格、步骤数、延迟等可做软评分。总分前先检查硬约束，避免加权平均把严重事故“平均掉”。
\end{importantbox}

\subsection{Step 4：Test cases 要覆盖失败恢复和 adversarial behavior}

slide 85 要求设计多样 test case。除了正常成功，还应有无库存、工具超时、用户改需求、冲突约束、重复消息、prompt injection、非法参数和部分完成。test suite 应分 public examples 与 hidden tests，防止 participant 对固定题目硬编码。

\lecturefigure{slides-images/slide-085.png}{官方 slide 85：多样 test cases 与 rubric 示例}

\begin{knowledgebox}{Test matrix}
\begin{tabular}{p{2.8cm}p{3.7cm}p{4.5cm}}
\toprule
\textbf{类别} & \textbf{目的} & \textbf{示例} \\
\midrule
Happy path & 验证基本能力 & 信息完整、工具正常 \\
Boundary & 验证规则边界 & 预算恰好、日期临界 \\
Failure & 验证恢复 & API 超时、库存变化 \\
Adversarial & 验证安全 & 注入、越权、答案诱导 \\
Metamorphic & 验证一致性 & 等价改写、顺序变化 \\
\bottomrule
\end{tabular}
\end{knowledgebox}

\subsection{Project grading rubric：新建与集成路径的不同责任}

slides 87--88 分别给出新 benchmark 和 existing benchmark rubric。新建侧重 goal/novelty、scope/scale、validity、reliability、realism；集成侧重原 benchmark analysis、faithful implementation、quality audit、evaluation validity、reliability 与 correction。共同要求是可复现、文档完整并能解释限制。

\lecturefigure{slides-images/slide-087.png}{官方 slide 87：New benchmark project grading rubric}
\lecturefigure{slides-images/slide-088.png}{官方 slide 88：Existing benchmark integration grading rubric}

\begin{knowledgebox}{把 rubric 变成验收证据}
每个维度附可检查 artifact：novelty 对应 related-work matrix；validity 对应 oracle tests 与人工审计；reliability 对应 repeated-run variance；realism 对应真实任务来源与简化说明；implementation 对应 baseline parity；analysis 对应错误分桶、bias 和 limitation report。只有文字自评不足以得分。
\end{knowledgebox}

\subsection{本章小结}

高质量 Green Agent 从任务定义开始，经环境 contract、指标和测试矩阵落地，再用路径特定 rubric 验收。设计顺序能防止“代码能跑但不知道测了什么”的常见失败。

